\chapter*{Abstract}
\addcontentsline{toc}{chapter}{Abstract}

Melanoma is survivable when found early and frequently fatal when it is not, so the
practical value of automated skin lesion classification lies in triage rather than
diagnosis. That places a constraint the accuracy literature rarely addresses: a triage
tool has to run on the device in front of the patient. ResNet-50 serialises to 94.3~MB
in the seven-class configuration used here; MobileNetV2 comes to 9.1~MB. The question is
what that reduction costs.

The published record answers it poorly. Studies comparing architectures on HAM10000 vary
the split protocol, augmentation, optimiser and schedule alongside the architecture, so a
reported gap between two backbones conflates the backbone with the setup. Headline
accuracy compounds the problem: HAM10000 is 66.9\% benign nevi, so a model that answers
nevus to every image scores 66.9\% without recognising a single melanoma.

This thesis runs the comparison as a controlled experiment. A single pipeline holds the
lesion-level split, augmentation, optimiser, batch size, schedule and hardware fixed, and
varies only the architecture and the class-imbalance strategy. The result is a
$2\times3$ factorial over MobileNetV2 and ResNet-50 crossed with augmentation-only,
weighted cross-entropy and random oversampling, extended by a six-run EfficientNet-B3
ablation at two input resolutions and analysed in a separate statistical family.
Performance is reported per class, with melanoma and basal cell carcinoma as the
headline, and comparisons are tested with McNemar's paired test under Holm-Bonferroni
correction.

Four findings emerge. MobileNetV2 achieved higher macro F1 than ResNet-50 under all
three strategies, and under two of them the architectures were statistically
indistinguishable on per-image correctness despite a ten-fold difference in parameter
count. The best imbalance strategy depended on the architecture: weighted cross-entropy
improved MobileNetV2's clinical recall while both remedies degraded ResNet-50's, which
Chapter~8 traces to optimisation stability rather than to class frequency. And selecting
checkpoints on validation loss, the protocol this work inherited, cost up to 0.067 macro
F1 and did so unevenly across architectures, which is a limitation of the design rather
than a property of the models and is reported as such.

\vspace{1em}
\noindent\textbf{Keywords:} skin lesion classification, HAM10000, class imbalance,
MobileNetV2, ResNet-50, EfficientNet-B3, controlled factorial design, McNemar's test,
deployable deep learning
